\section*{Capítulo \ref{ch:b2:eruptive-styles} --- As erupções vulcânicas}

\begin{solution}{exo:b2:eruptive-styles:1}
(a) $\qty{1.3}{km^3} = \qty{1.3e9}{m^3}$, entre $10^9$ e $10^{10}$: VEI~5.
(b) $\qty{2.9e8}{m^3}$: VEI~4. (c) $\qty{2e6}{m^3}$, entre $10^6$ e $10^7$:
VEI~2. (d) O VEI~6 começa em \qty{e10}{m^3}, de modo que a \omterm{def:b2:eruptive-styles:tephra}{tefra} deve ter um volume aparente de pelo menos
\qty{10}{km^3}; ``mais de \qty{5}{km^3}'' é um limite inferior, não uma contradição.
\end{solution}

\begin{solution}{exo:b2:eruptive-styles:2}
(a) Cinza (mais fina que \qty{2}{mm}). (b) Um lápilo (\qtyrange{2}{64}{mm}).
(c) Um bloco: maior que \qty{64}{mm}, anguloso, sólido ao ser ejetado. (d) Uma
\omterm{def:b2:eruptive-styles:tephra}{bomba vulcânica}, uma bomba fusiforme: maior que \qty{64}{mm} e moldada enquanto fundida,
em voo.
\end{solution}

\begin{solution}{exo:b2:eruptive-styles:3}
(a) $0.304\times(10^6)^{0.241} = 0.304\times27.9 = \qty{8.5}{km}$;
$0.304\times(10^8)^{0.241} = 0.304\times84.7 = \qty{25.8}{km}$.
(b) $Q = (35/0.304)^{1/0.241} = 115.1^{4.149} = \qty{3.6e8}{kg/s}$, e $I =
8.55 + 3 = 11.6$. (c) $Q \propto H^{4.15}$, de modo que uma altura \qty{20}{\percent} alta demais
superestima $Q$ por um fator $1.2^{4.15} = 2.1$.
\end{solution}

\begin{solution}{exo:b2:eruptive-styles:4}
$V_{\mathrm{DRE}} = 0.29\times1000/2500 = \qty{0.116}{km^3}$; $m =
\num{0.29e9}\times1000 = \qty{2.9e11}{kg}$; $M = \log_{10}(\num{2.9e11}) - 7 = 11.46
- 7 = 4.46$.
\end{solution}

\begin{solution}{exo:b2:eruptive-styles:5}
$k = \ln(40/10)/(15 - 5) = 1.386/10 = \qty{0.139}{km^{-1}}$; $T_0 = 40\,e^{0.139\times5}
= 40\times2.00 = \qty{80}{cm}$; $\ln2/k = \qty{5.0}{km}$. Com $T_0 =
\qty{8.0e-4}{km}$: $V = 2\times\num{8.0e-4}/0.139^2 = \qty{0.083}{km^3} =
\qty{8.3e7}{m^3}$, VEI~3.
\end{solution}

\begin{solution}{exo:b2:eruptive-styles:6}
(a) $\int_0^x T_0e^{-ky}2y\,dy = 2T_0\big[-\tfrac{y}{k}e^{-ky} - \tfrac{1}{k^2}e^{-ky}\big]_0^x
= \tfrac{2T_0}{k^2}[1 - (1 + kx)e^{-kx}]$. (b) Iterando $u = \ln(2(1 + u))$ a partir de
$u = 1$: $1.386$, $1.564$, $1.635$, $1.662$, $1.672$, \dots, que converge para
$u = 1.678$, para o qual $(1 + u)e^{-u} = 2.678\times0.1867 = 0.500$. (c) $x_{1/2} =
1.678/0.139 = \qty{12.1}{km}$, dentro da isópaca mais externa ($\sqrt A =
\qty{15}{km}$); no entanto, o volume além dessa isópaca é $(1 + 2.09)e^{-2.09} =
\qty{38}{\percent}$ do total. Quase dois quintos do depósito ficam onde
nada foi medido: o volume repousa sobre a extrapolação.
\end{solution}

\begin{solution}{exo:b2:eruptive-styles:7}
(a) $p_s = (5.0/0.444)^2 = \qty{127}{MPa}$, a $127/0.025 = \qty{5.1}{km}$.
(b) $x = (5.0 - 0.444\sqrt{10})/100 = (5.0 - 1.404)/100 = 0.0360$;
$V_g = 0.0360\times8.314\times1123/(0.018015\times10^7) = \qty{1.86e-3}{m^3}$;
$V_m = 0.964/2300 = \qty{4.19e-4}{m^3}$; $\phi = 1.86/(1.86 + 0.419) = 0.82$.
(c) Um fundido basáltico é de milhares a bilhões de vezes menos viscoso que um
fundido riolítico (\cref{ch:b2:magma-emplacement}): suas bolhas sobem,
coalescem e escapam, e o gás sai por uma espuma permeável ou pela
boca, de modo que o \omterm{def:b2:magma-emplacement:magma}{magma} se desgaseifica em sistema aberto em vez de reter seu gás; sua
fragmentação, quando ocorre, fica restrita a uma fonte perto da boca. Uma ascensão mais
lenta deixa mais tempo para isso.
\end{solution}

\begin{solution}{exo:b2:eruptive-styles:8}
(a) $u_z(0) = 0.75\times10^7/(\pi\times(5000)^2) = \qty{0.095}{m} = \qty{95}{mm}$.
(b) $z_0 = 3.8/0.766 = \qty{5.0}{km}$. (c) $(1 + 1)^{-3/2} = 0.35$ do
máximo.
\end{solution}

\begin{solution}{exo:b2:eruptive-styles:9}
O maior alcance é a \ang{45}: $R = v^2/g = 100^2/9.81 = \qty{1.0}{km}$, e
$200^2/9.81 = \qty{4.1}{km}$. A resistência do ar encurta o voo, de modo que esses são
limites superiores, próximos da verdade para blocos grandes e longos demais para os pequenos.
Como $R = v^2/g$, a zona de perigo cresce com o quadrado da velocidade de lançamento:
dobrar a velocidade da explosão quadruplica seu raio.
\end{solution}

\begin{solution}{exo:b2:eruptive-styles:10}
(a) Com $T = T_0e^{-k_1x}$ para $x < X$ e $T = T_Xe^{-k_2(x - X)}$ além, $T_X =
T_0e^{-k_1X}$ (continuidade), use-se $\int_0^X ye^{-ky}dy = \frac{1}{k^2} - e^{-kX}\big(\frac{X}{k} +
\frac{1}{k^2}\big)$ e $\int_X^\infty ye^{-ky}dy = e^{-kX}\big(\frac{X}{k} +
\frac{1}{k^2}\big)$:
a parte proximal dá $2T_0\big[\frac{1}{k_1^2} - e^{-k_1X}\frac{k_1X + 1}{k_1^2}\big]$ e a
parte distal $2T_Xe^{k_2X}\int_X^\infty ye^{-k_2y}dy = 2T_X\frac{k_2X + 1}{k_2^2}$; com $T_X =
T_0e^{-k_1X}$, sua soma é a fórmula do quadro de método.
(b) $T_0 = \qty{0.01285}{km}$: $2T_0/k_1^2 = 0.0257/0.1332 = \qty{0.193}{km^3}$;
$e^{-0.365\times8.56} = 0.0440$; $(0.138\times8.56 + 1)/0.138^2 = 114.5$ e
$(0.365\times8.56 + 1)/0.365^2 = 30.96$; a correção é $0.0257\times0.0440\times83.6
= \qty{0.094}{km^3}$ e $V = \qty{0.287}{km^3}$, os \qty{0.29}{km^3} publicados.
Além de $X$ ficam $2T_X(X/k_2 + 1/k_2^2) = 2\times\num{5.65e-4}\times114.5 =
\qty{0.129}{km^3}$, isto é, \qty{45}{\percent}.
\end{solution}

\begin{solution}{exo:b2:eruptive-styles:11}
(a) $100^{1/4} = 3.16$; as leis empíricas dão $100^{0.241} = 3.03$ e
$100^{0.226} = 2.83$. (b) $H \propto N^{-3/4}$: uma atmosfera mais fortemente estratificada
detém a coluna mais abaixo; para $N$ maior em \qty{20}{\percent},
$1.2^{-0.75} = 0.87$, uma altura \qty{13}{\percent} menor. (c) As plumas fracas são
inclinadas pelo vento, o que as rebaixa, enquanto o ar úmido libera nelas calor latente,
o que as eleva; esses efeitos pesam de modo diferente sobre as pequenas e as grandes
erupções e não estão na lei simples. E as vazões só são conhecidas a menos de
fatores de vários, um erro na variável do eixo horizontal que enviesa
uma inclinação ajustada em direção a zero.
\end{solution}

\begin{solution}{exo:b2:eruptive-styles:12}
(a) $132/400\times100 = 33$ por século antes de 1900; $95/125\times100 = 76$ por
século desde então. (b) $33/76 = \qty{43}{\percent}$. (c) $21/56 = \qty{38}{\percent}$:
o mesmo sub-registro, por um fator de dois a três. (d) As erupções maiores são
vistas de mais longe e deixam depósitos mais espessos e extensos que resistem à erosão
por mais tempo; uma erupção de VEI~4 em uma ilha remota em 1700 podia passar sem registro, uma
erupção de VEI~6 dificilmente.
\end{solution}

\begin{solution}{pb:b2:eruptive-styles:1}
\textbf{1.} $\ln T = 5.298$, $4.605$, $3.871$, $2.639$, $1.386$. As inclinações entre
isópacas sucessivas são $-0.0693$, $-0.0367$, $-0.0308$ e $-0.0313$ por km: um
segmento íngreme entre 10 e \qty{20}{km}, um suave de 40 a
\qty{120}{km}; a quebra fica entre 20 e \qty{40}{km}.

\textbf{2.} $k_1 = \ln(200/100)/10 = \qty{0.0693}{km^{-1}}$; $T_0 = 200\,e^{0.693} =
\qty{400}{cm}$.

\textbf{3.} $167e^{-1.244} = \qty{48.1}{cm}$, $167e^{-2.488} = \qty{13.9}{cm}$,
$167e^{-3.732} = \qty{4.00}{cm}$: a menos de \qty{1}{\percent} de 48, 14 e 4.0.

\textbf{4.} $400e^{-0.0693X} = 167e^{-0.0311X}$ dá $X = \ln(400/167)/(0.0693 -
0.0311) = 0.873/0.0382 = \qty{22.9}{km}$.

\textbf{5.} $\ln2/0.0693 = \qty{10.0}{km}$ e $\ln2/0.0311 = \qty{22.3}{km}$.

\textbf{6.} $T_0 = \qty{0.0040}{km}$: $V = 2\times0.0040/0.0693^2 =
\qty{1.67}{km^3}$.

\textbf{7.} $e^{-0.0693\times22.9} = 0.205$; $(0.0311\times22.9 + 1)/0.0311^2 = 1770$;
$(0.0693\times22.9 + 1)/0.0693^2 = 539$; correção $2\times0.0040\times0.205\times1231 =
\qty{2.02}{km^3}$; $V = 1.67 + 2.02 = \qty{3.68}{km^3}$, dos quais a correção representa
\qty{55}{\percent}.

\textbf{8.} $3.68/1.67 = 2.2$. As cinzas distais finas são delgadas, facilmente erodidas,
retrabalhadas ou caídas no mar; em um depósito antigo, muitas vezes estão perdidas ou nunca foram
cartografadas, e o volume é então subestimado por um fator de dois ou mais.

\textbf{9.} $\qty{3.68e9}{m^3}$ fica entre $10^9$ e $10^{10}$: VEI~5.

\textbf{10.} $V_{\mathrm{DRE}} = 3.68\times900/2400 = \qty{1.38}{km^3}$.

\textbf{11.} $m = \num{3.68e9}\times900 = \qty{3.31e12}{kg}$.

\textbf{12.} $M = \log_{10}(\num{3.31e12}) - 7 = 12.52 - 7 = 5.52$.

\textbf{13.} $3.68/1.3 = 2.8$ vezes.

\textbf{14.} Uma coluna mais alta que \qty{25}{km}, e um estilo pliniano.

\textbf{15.} $Q = \num{3.31e12}/36000 = \qty{9.2e7}{kg/s}$; $I = 7.96 + 3 = 11.0$.

\textbf{16.} $0.304\times(\num{9.2e7})^{0.241} = 0.304\times83.0 = \qty{25.2}{km}$.

\textbf{17.} $0.345\times(\num{9.2e7})^{0.226} = 0.345\times63.1 = \qty{21.8}{km}$.

\textbf{18.} Os \qty{31}{km} medidos ficam \qty{23}{\percent} acima da primeira
previsão e \qty{42}{\percent} acima da segunda, dentro da dispersão de
\qtyrange{53}{60}{\percent} das leis: compatível.

\textbf{19.} $Q = (31/0.304)^{1/0.241} = 102.0^{4.149} = \qty{2.2e8}{kg/s}$, 2.3 vezes
a vazão média. Os clastos registram a coluna mais alta, no pico da
erupção, enquanto a questão~15 faz a média sobre toda a fase; e parte da massa
pode estar ausente do depósito (cinzas distais, \omterm{def:b2:eruptive-styles:pdc}{fluxos piroclásticos}).

\textbf{20.} Sim: de 22 a \qty{31}{km}, bem acima da tropopausa, a \qty{16}{km}.
O gás levado à estratosfera se espalha em torno do globo; o gás-chave é o
dióxido de enxofre, e os vinte milhões de toneladas que o Pinatubo lá injetou em 1991
resfriaram o globo em cerca de meio grau durante dois anos.

\textbf{21.} $\lambda_3 = 399/75 = 5.32$, $\lambda_4 = 95/125 = 0.760$, $\lambda_5 = 22/225 =
0.0978$, $\lambda_6 = 8/525 = 0.0152$ por ano; logaritmos $0.726$, $-0.119$, $-1.010$,
$-1.817$.

\textbf{22.} A reta dá $10^{0.73} = 5.37$, $10^{-0.12} = 0.759$, $10^{-0.97} =
0.107$ e $10^{-1.82} = 0.0151$: a \qty{1}{\percent}, \qty{0}{\percent},
\qty{10}{\percent} e \qty{1}{\percent} das contagens.

\textbf{23.} O registro das grandes erupções permanece completo mais para trás no tempo,
porque elas são notadas de mais longe e seus depósitos sobrevivem por mais tempo; uma
janela que começasse em 1500 subestimaria muito as erupções de VEI~3, enquanto uma janela
que começasse em 1950 não conteria quase nenhuma erupção de VEI~6.

\textbf{24.} $\lambda_5 = \num{0.107}$ por ano, um período de retorno de $1/0.107 =
\num{9.3}\text{ anos}$; $1 - e^{-1.07} = 1 - 0.343 = \qty{66}{\percent}$ em dez anos.

\textbf{25.} $\lambda_7 = 10^{3.28 - 5.95} = 10^{-2.67} = 0.00214$ por ano, uma a cada
470 anos; o registro dá $2/1025 = 0.00195$ por ano: concordância a menos de
\qty{10}{\percent}, para um ponto não usado no ajuste.

\textbf{26.} $\lambda_6 = 10^{3.28 - 5.10} = 10^{-1.82} = 0.0151$ por ano; $\lambda t = 1.51$
e $1 - e^{-1.51} = 1 - 0.220 = 0.780$: a probabilidade de pelo menos uma erupção
de VEI~6 ou mais em algum lugar da Terra nos próximos cem anos é de
\textbf{\qty[text-series-to-math=true]{78}{\percent}}. Ela repousa sobre uma taxa constante,
sobre erupções independentes umas das outras (um processo de Poisson) e sobre um
catálogo completo dentro de suas janelas.
\end{solution}
